\documentclass[pdflatex,sn-mathphys-num]{sn-jnl}% Math and Physical Sciences Numbered Reference Style
%%\documentclass[pdflatex,sn-nature]{sn-jnl}% Nature Portfolio
%%\documentclass[pdflatex,sn-basic]{sn-jnl}% Basic Springer Nature / Chemistry
%%\documentclass[pdflatex,sn-aps]{sn-jnl}% American Physical Society
%%\documentclass[pdflatex,sn-vancouver-num]{sn-jnl}% Vancouver Numbered

\usepackage{graphicx}%
\usepackage{multirow}%
\usepackage{amsmath,amssymb,amsfonts}%
\usepackage{amsthm}%
\usepackage{mathrsfs}%
\usepackage[title]{appendix}%
\usepackage{xcolor}%
\usepackage{textcomp}%
\usepackage{manyfoot}%
\usepackage{booktabs}%

\theoremstyle{thmstyleone}%
\newtheorem{theorem}{Theorem}%
\newtheorem{proposition}[theorem]{Proposition}%

\theoremstyle{thmstyletwo}%
\newtheorem{example}{Example}%
\newtheorem{remark}{Remark}%

\theoremstyle{thmstylethree}%
\newtheorem{definition}{Definition}%

\raggedbottom

\begin{document}

\title[]{Practical prediction of the production output of multi-isotope generators}

\author*[1]{\fnm{Artem} \sur{Lebedev}}\email{artem@indiechemistry.com}

\affil*[1]{\orgdiv{Department of Chemistry and Chemical Biology}, \orgname{McMaster University}, \orgaddress{\city{Hamilton, ON}, \country{Canada}}}

\abstract{A brief explanation of the math and code behind practical analysis of complex isotope generators.}

\maketitle

\section{Introduction}
\label{sec:introduction}

Isotope generators are the backbone of medical isotope production, accounting for well over 90\% of all the nuclear medicine procedures. The bulk of this volume comes from the use of Mo-99/Tc-99m generator, but other parent-daughter pairs have been gaining popularity recently, namely Ge-68/Ga-68, Ra-225/Ac-225 and a few others.

Theoretical calculations of the generator state and predictions of the generator output rely on basic equations suggested by Rutherford and Soddy \cite{rutherford_lx_1903}. Interestingly enough, already in their foundational paper from 1903 they mention multi nuclide decay chain that is the foundation of the ultra-modern Pb-212 generators.

Few years later, Bateman generalized equations from \cite{rutherford_lx_1903} into a system of differential equations and suggested analytical solution now known as Bateman equations \cite{bateman_solution_1908}. These equations, allow to predict the present state of the generator, that is the concentration of all isotopes, given the initial state of the generator. While this is technically sufficient, practical considerations of harvesting schedules, generator shelf-life and economical viability add a new layer of complexity. Additionally, computational implementations of Bateman equations are known to suffer from a series of drawbacks \cite{bakin_computational_2018}.

This paper describes the logic behind a computer code designed to simulate generator states and output under various operational conditions. 


\section{Two solutions to the same equations}
\label{sec:two_solutions}

Let's consider the most important isotope generator in nuclear medicine: Tc generator, that is still a staple of the clinical practice, used for millions upon millions of SPECT scans every year, all around the world. Mo-99 decays two ways: \~{}87.6\% to Tc-99m (branching ratio $b$) and the remaining \~{}12.4\% directly to (virtually) stable Tc-99, which does not factor into the generator considerations. From here, a conventional chemist's thinking can easily lead us to the following system of equations:

\begin{equation}
	\begin{cases}
		\dfrac{d[^{99}Mo]}{dt} &= -\lambda_{Mo}[^{99}Mo] \\ \\
		\dfrac{d[^{99m}Tc]}{dt} &= b\lambda_{Mo}[^{99}Mo]-\lambda_{Tc}[^{99m}Tc]
	\end{cases}
\end{equation}

where $\lambda_{Mo} = \ln2/65.94h$, $\lambda_{Tc} = \ln2/6.01h$ and branching ratio $b = 0.876$

What this is saying is that the rate of accumulation of \textsuperscript{99}Mo is negatively proportional to its concentration, i.e. the more \textsuperscript{99}Mo we have in solution the slower it accumulates, or what is the same, the faster it decays. This is essentially a definition of the radioactive decay. For the \textsuperscript{99m}Tc the picture is more complex, but is still intuitive: its rate of accumulation is \textit{negatively} proportional to its own concentration $[{^{99m}Tc}]$, but also \textit{positively} proportional to the parent's concentration $[^{99}Mo]$. Also, let's not forget that the $[^{99}Mo]$ is scaled to the branching ratio in the second equation. 

**Matrix form**

\$\$A = \textbackslash{}begin\{pmatrix\} -\textbackslash{}lambda\_1 \& 0 \textbackslash{}\textbackslash{} b\textbackslash{}lambda\_1 \& -\textbackslash{}lambda\_2 \textbackslash{}end\{pmatrix\}\$\$

Only the subdiagonal entry picks up the factor \$b\$. Diagonalizing exactly as before (eigenvalues still \$-\textbackslash{}lambda\_1, -\textbackslash{}lambda\_2\$ since those live on the diagonal regardless of \$b\$):

\$\$e\^{}\{At\} = \textbackslash{}begin\{pmatrix\} e\^{}\{-\textbackslash{}lambda\_1 t\} \& 0 \textbackslash{}\textbackslash{}[4pt] \textbackslash{}dfrac\{b\textbackslash{}lambda\_1\}\{\textbackslash{}lambda\_2-\textbackslash{}lambda\_1\}\textbackslash{}left(e\^{}\{-\textbackslash{}lambda\_1 t\}-e\^{}\{-\textbackslash{}lambda\_2 t\}\textbackslash{}right) \& e\^{}\{-\textbackslash{}lambda\_2 t\} \textbackslash{}end\{pmatrix\}\$\$

So Bateman becomes:

\$\$N\_2(t) = N\_1\^{}0\textbackslash{},\textbackslash{}frac\{b\textbackslash{}lambda\_1\}\{\textbackslash{}lambda\_2-\textbackslash{}lambda\_1\}\textbackslash{}left(e\^{}\{-\textbackslash{}lambda\_1 t\}-e\^{}\{-\textbackslash{}lambda\_2 t\}\textbackslash{}right)\$\$

— literally just the un-branched formula times \$b\$. This makes sense: branching ratio only ever multiplies the *production* term, never the decay-loss term, so it factors out linearly and leaves the transient/secular equilibrium math (the classic 6-hour Mo-99/Tc-99m generator ``milking'' curve) otherwise unchanged.

**Where it stops being this simple**

If you extend the chain further (e.g. tracking Tc-99 ground state buildup too, or handling isotopes with 3+ branches like your Ac-225 chain has at Bi-213 → Tl-209/Po-213), you get multiple daughters fed from one parent. That's not a single chain anymore — \$A\$ becomes a genuine tree structure, and you'd want one column feeding multiple rows (each weighted by its own \$b\_i\$, with \$\textbackslash{}sum b\_i = 1\$) rather than a simple bidiagonal matrix. The matrix exponential machinery still works fine, it's just no longer tridiagonal/bidiagonal, so `expm` is basically mandatory rather than a convenience at that point.

\backmatter

\bmhead{Acknowledgements}
Author thanks Ian Horwath for a stimulating discussion and Andrew Knight for the feedback on the web app version of these computations.

\section*{Declarations}

\begin{itemize}
\item Funding: The work was entirely self-funded
\item Conflict of interest: The author is an employee of BWXT Medical, Ltd.
\item Data availability: The article and the app code are available on GitHub
\end{itemize}

\bibliography{bibliography/bateman_rev.bib}

\end{document}
